\chapter{Natural numbers}
\label{ch:naturals}

We now construct the numbers $0, 1, 2, 3, \dots$ out of sets. The idea, due to
von Neumann, is to let \emph{each number be the set of all numbers before it}:
\[
  0 = \varnothing, \quad 1 = \set{0}, \quad 2 = \set{0, 1}, \quad 3 = \set{0, 1,
  2}, \quad \dots
\]
Then the number $n$ has exactly $n$ elements, \enquote{$m$ is smaller than $n$}
can simply mean $m \in n$, and passing to the next number means adding the number
itself as a new element.

Then we prove that these sets behave the way numbers should. The \emph{Peano
axioms}, induction included, come out as \emph{theorems}. So do addition,
multiplication and order, with all the usual laws.

\section{Construction}

\begin{definition}[Successor, inductive set]
\label{def:inductive}
The \emph{successor} of a set $x$ is the set $x \cup \set{x}$. A set $I$ is
\emph{inductive} if $\varnothing \in I$ and $\forall x\; (x \in I \Rightarrow x
\cup \set{x} \in I)$.
\end{definition}

An inductive set contains $\varnothing$, $\varnothing \cup \set{\varnothing}$, the
successor of that, and so on. None of the axioms so far guarantees that such a set
exists: every set we have built so far is finite. So we need a new axiom.

\begin{axiom}[Infinity]
\label{ax:infinity}
There is an inductive set.
\end{axiom}

\begin{definition}[Natural numbers]
\label{def:naturals}
Let $I_0$ be an inductive set (\cref{ax:infinity}) and
\[
  \N \coloneqq \setof{x \in I_0}{x \in I \text{ for every inductive set } I}.
\]
The elements of $\N$ are the \emph{natural numbers}. The \emph{numerals} are
\[
  0 \coloneqq \varnothing,\quad 1 \coloneqq 0 \cup \set{0},\quad 2 \coloneqq 1 \cup
  \set{1},\quad \dots,\quad 9 \coloneqq 8 \cup \set{8}.
\]
\register{naturals}{$\N$}{the set of natural numbers}
\register{naturals}{$0, 1, \dots, 9$}{the first ten natural numbers; $0$ is the
empty set regarded as a number (\cref{rule:identification})}
\end{definition}

$\N$ doesn't depend on which $I_0$ we picked: if $I_1$ is another inductive set,
any $x$ in every inductive set is in both $I_0$ and $I_1$, so both choices give
the same elements. Two facts follow directly from the definition:
\begin{itemize}
  \item \emph{$\N$ is inductive.} $\varnothing$ lies in every inductive set. If
    $x$ lies in every inductive set, so does $x \cup \set{x}$.
  \item \emph{$\N$ is contained in every inductive set.} That is how its elements
    were chosen.
\end{itemize}
So $\N$ is the \emph{smallest} inductive set. Indeed $1 = \varnothing \cup
\set{\varnothing} = \set{0}$, $2 = \set{0} \cup \set{1} = \set{0, 1}$, and so on.
Positional notation like $10, 11, \dots$ needs division with remainder and waits
for the chapter on divisibility. In this book, $0 \in \N$.

\section{Induction}

\begin{theorem}[Principle of induction]
\label{thm:induction}
Let $X \subseteq \N$ with
\begin{enumerate}
  \item $0 \in X$, and
  \item $n \in X \Rightarrow n \cup \set{n} \in X$ for all $n \in \N$.
\end{enumerate}
Then $X = \N$.

Equivalently, for a predicate $P$: if $P(0)$ holds (the \emph{base case}) and $P(n)
\Rightarrow P(n \cup \set{n})$ holds for every $n \in \N$ (the \emph{induction
step}, with \emph{induction hypothesis} $P(n)$), then $P(n)$ holds for all $n \in
\N$.
\end{theorem}

\begin{proof}
By (1) and (2), $X$ is inductive: every element of $X$ lies in $\N$, so (2)
applies to it. So $\N \subseteq X$ (\cref{def:naturals}), and $X \subseteq \N$ is
given. For the predicate form, apply this to $X \coloneqq \setof{n \in \N}{P(n)}$.
\end{proof}

Induction is not a new law of logic. It follows from the way $\N$ was built: as
the \emph{smallest} inductive set, it has no room for anything that isn't reached
from $0$ by successors.

\begin{lemma}[Shape of natural numbers]
\label{lem:shape}
For every $n \in \N$:
\begin{enumerate}
  \item $n \subseteq \N$: every element of a natural number is a natural number;
  \item $n$ is \emph{transitive}: $m \in n \Rightarrow m \subseteq n$;
  \item $n \notin n$.
\end{enumerate}
\end{lemma}

\begin{proof}
Each by induction on $n$.

(1) $0 = \varnothing \subseteq \N$. If $n \subseteq \N$, then $n \cup \set{n} \subseteq \N$,
because $n \in \N$.

(2) $0$ has no elements, so the condition holds vacuously. Assume $n$ is
transitive and let $m \in n \cup \set{n}$. Then $m \in n$ or $m = n$. If $m \in n$, then $m
\subseteq n \subseteq n \cup \set{n}$. If $m = n$, then $m = n \subseteq n \cup \set{n}$.

(3) $0 \notin 0$, since $0$ has no elements. Assume $n \notin n$, and suppose $n \cup \set{n}
\in n \cup \set{n}$. Then $n \cup \set{n} \in n$ or $n \cup \set{n} = n$. If $n \cup \set{n} = n$, then $n \in n \cup \set{n} = n$. If
$n \cup \set{n} \in n$, then $n \cup \set{n} \subseteq n$ by (2), and $n \in n \cup \set{n}$ gives $n \in n$. Both
contradict the induction hypothesis.
\end{proof}

\begin{definition}[Successor map]
\label{def:successor-map}
Since $\N$ is inductive, $n \cup \set{n} \in \N$ for every $n \in \N$. The
\emph{successor map} is the map $S \colon \N \to \N$ with $n \via{S} n \cup
\set{n}$.
\register{naturals}{$S$}{the successor map $\N \to \N$, sending $n$ to $n + 1$}
\end{definition}

\begin{theorem}[Peano axioms]
\label{thm:peano}
\leavevmode
\begin{enumerate}
  \item $0 \in \N$.
  \item $S$ is a map $\N \to \N$.
  \item $0$ is not a value of $S$: $n \cup \set{n} \neq 0$ for every $n \in \N$.
  \item $S$ is injective: $m \cup \set{m} = n \cup \set{n} \Rightarrow m = n$.
  \item Induction (\cref{thm:induction}).
\end{enumerate}
\end{theorem}

\begin{proof}
(1) and (2): $\N$ is inductive. (3): $n \in n \cup \set{n}$, so $n \cup \set{n}
\neq \varnothing$.

(4) Let $m \cup \set{m} = n \cup \set{n}$ and suppose $m \neq n$. Since $m$ is an
element of the left side, $m \in n$ or $m = n$, so $m \in n$. In the same way $n
\in m$. By transitivity (\cref{lem:shape}(2)), $m \in n$ gives $m \subseteq n$,
and with $n \in m$ that gives $n \in n$, contradicting \cref{lem:shape}(3).
\end{proof}

\begin{corollary}[Predecessor]
\label{cor:predecessor}
For every $n \in \N$ with $n \neq 0$ there is exactly one $p \in \N$ with $n = p
\cup \set{p}$, the \emph{predecessor} of $n$.
\end{corollary}

\begin{proof}
Existence, by induction on the statement \enquote{$n = 0$ or $n = p \cup \set{p}$
for some $p \in \N$}: it holds for $0$, and for $n \cup \set{n}$ with $p \coloneqq
n$. Uniqueness: \cref{thm:peano}(4).
\end{proof}

\section{Recursion}

Induction \emph{proves} statements about all $n$. Recursion \emph{defines} maps on
$\N$: give the value at $0$, and a rule that turns the value at $n$ into the value
at the successor of $n$. That this defines a map at all is a theorem.

\begin{theorem}[Recursion theorem]
\label{thm:recursion}
Let $A$ be a set, $a \in A$, and $g \colon A \to A$ a map. Then there is exactly
one map $h \colon \N \to A$ with
\[
  0 \via{h} a \qquad\text{and}\qquad n \via{h} y \via{g} z \;\Rightarrow\; n \cup
  \set{n} \via{h} z \quad \text{for all } n \in \N,\; y, z \in A.
\]
In terms of composition, the second condition says $h \circ S = g \circ h$.
\end{theorem}

\begin{proof}
\emph{Idea:} $h$ should be the smallest relation that relates $0$ to $a$ and is
closed under the step from $n$ to its successor. Then we show this relation is a
map.

Call a relation $R \subseteq \N \times A$ \emph{closed} if $0 \mathrel{R} a$ and $n
\mathrel{R} y \land y \via{g} z \Rightarrow (n \cup \set{n}) \mathrel{R} z$. Then
$\N \times A$ is closed. Let $h$ be the intersection of all closed relations,
i.e.\ of the nonempty set of closed elements of $\Pow (\N \times A)$
(\cref{def:set-operations}). Then $h$ is closed as well, and $h \subseteq R$ for
every closed $R$.

\emph{$h$ is a map.} We show by induction on $n$: there is exactly one $y$ with $n
\mathrel{h} y$.
\begin{itemize}
  \item $n = 0$. We have $0 \mathrel{h} a$. Suppose $0 \mathrel{h} b$ with $b \neq
    a$. Then $h' \coloneqq h \setminus \set{(0, b)}$ is still closed: it relates $0$
    to $a$, and every pair added by the step has as first component a successor,
    which is not $0$ (\cref{thm:peano}(3)). That contradicts $h \subseteq h'$,
    since $(0, b) \notin h'$.
  \item Step. Assume $y$ is the only element with $n \mathrel{h} y$, and let $z$ be
    the value of $g$ at $y$. Since $h$ is closed, $(n \cup \set{n}) \mathrel{h} z$.
    Suppose $(n \cup \set{n}) \mathrel{h} w$ with $w \neq z$. Then $h' \coloneqq h
    \setminus \set{(n \cup \set{n}, w)}$ is still closed. It relates $0$ to $a$.
    If $k \mathrel{h'} u \via{g} v$ and $(k \cup \set{k}, v) = (n \cup \set{n},
    w)$, then $k = n$ (\cref{thm:peano}(4)), so $u = y$ by the induction
    hypothesis, so $v = z$, which contradicts $v = w \neq z$. Again $h \subseteq h'$
    fails.
\end{itemize}
So $h$ is a map, and since it is closed, it has both properties.

\emph{Uniqueness.} If $h'$ also has both properties, we show by induction that $h$
and $h'$ have the same value at every $n$: both send $0$ to $a$, and if both send
$n$ to $y$, and $y \via{g} z$, both send $n \cup \set{n}$ to $z$. Use
\cref{thm:map-equality}.
\end{proof}

Often the step also needs to know $n$, as for the factorial. The stronger form
follows from the simple one.

\begin{corollary}[Recursion with index]
\label{cor:recursion-index}
Let $a \in A$ and $g \colon \N \times A \to A$. Then there is exactly one $h \colon
\N \to A$ with
\[
  0 \via{h} a \qquad\text{and}\qquad n \via{h} y \land (n, y) \via{g} z
  \;\Rightarrow\; n \cup \set{n} \via{h} z.
\]
\end{corollary}

\begin{proof}
Let $G \colon \N \times A \to \N \times A$ be the map that sends $(n, y)$ to $(n
\cup \set{n}, z)$, where $z$ is the value of $g$ at $(n, y)$ (\cref{thm:map-rule}).
\Cref{thm:recursion}, applied to $\N \times A$, the start value $(0, a)$ and $G$,
gives a map $H \colon \N \to \N \times A$. By induction, for every $n$ there is a
$y$ with $n \via{H} (n, y)$: this holds for $0$, and if $n \via{H} (n, y) \via{G} (n
\cup \set{n}, z)$, then $n \cup \set{n} \via{H} (n \cup \set{n}, z)$. So
\[
  h \coloneqq \setof{(n, y) \in \N \times A}{n \via{H} (n, y)}
\]
is a map with the two properties. Uniqueness again by induction.
\end{proof}

\section{Addition}

\begin{definition}[Addition]
\label{def:addition}
Let $m \in \N$. By \cref{thm:recursion} (with $A = \N$, $a = m$ and $g = S$) there
is exactly one map $h \colon \N \to \N$ with $0 \via{h} m$ and $h \circ S = S
\circ h$. \emph{Addition} is the operation on $\N$ (\cref{def:operation}) given by
\[
  (m, n) \via{+} k \;\defiff\; n \via{h} k \quad \text{for this } h,
\]
which exists by \cref{thm:map-rule}. We write $m + n$ for its value. It is
characterised by
\[
  \text{(A1)}\quad m + 0 = m, \qquad\qquad \text{(A2)}\quad m + (n \cup \set{n}) =
  (m + n) \cup \set{m + n}.
\]
\register{naturals}{$m + n$}{sum of $m$ and $n$}
\end{definition}

Since $1 = 0 \cup \set{0}$, (A1) and (A2) give $n + 1 = (n + 0) \cup \set{n + 0} =
n \cup \set{n}$. So the successor of $n$ is $n + 1$, and by
\cref{rule:abbreviation} \emph{from now on only $n + 1$ is written}. In particular
$n \via{S} n + 1$, and (A2) reads $m + (n + 1) = (m + n) + 1$. For example, since
$2 = 1 + 1$,
\[
  2 + 2 = 2 + (1 + 1) = (2 + 1) + 1 = 3 + 1 = 4.
\]

\begin{theorem}[Laws of addition]
\label{thm:addition-laws}
For all $k, m, n \in \N$:
\begin{enumerate}
  \item $0 + n = n$;
  \item $(m + 1) + n = (m + n) + 1$;
  \item $m + n = n + m$ (commutativity);
  \item $(k + m) + n = k + (m + n)$ (associativity);
  \item $k + n = m + n \Rightarrow k = m$ (cancellation);
  \item $m + n = 0 \Rightarrow m = 0 \land n = 0$.
\end{enumerate}
\end{theorem}

\begin{proof}
(1)--(5) by induction on $n$, using (A1) and (A2) at each step.

(1) $0 + 0 = 0$. If $0 + n = n$, then $0 + (n + 1) = (0 + n) + 1 = n + 1$.

(2) $(m + 1) + 0 = m + 1 = (m + 0) + 1$. If it holds for $n$, then $(m + 1) + (n +
1) = ((m + 1) + n) + 1 = ((m + n) + 1) + 1 = (m + (n + 1)) + 1$.

(3) $m + 0 = m = 0 + m$ by (1). If $m + n = n + m$, then $m + (n + 1) = (m + n) + 1
= (n + m) + 1 = (n + 1) + m$ by (2).

(4) $(k + m) + 0 = k + m = k + (m + 0)$. If it holds for $n$, then $(k + m) + (n +
1) = ((k + m) + n) + 1 = (k + (m + n)) + 1 = k + ((m + n) + 1) = k + (m + (n +
1))$.

(5) For $n = 0$ use (A1). Step: $k + (n + 1) = m + (n + 1)$ means $(k + n) + 1 = (m
+ n) + 1$, so $k + n = m + n$ by \cref{thm:peano}(4), so $k = m$ by the induction
hypothesis.

(6) If $n \neq 0$, then $n = p + 1$ (\cref{cor:predecessor}), and $m + n = (m + p)
+ 1 \neq 0$ (\cref{thm:peano}(3)). So $n = 0$, and then $m = m + 0 = 0$.
\end{proof}

By associativity, brackets in sums like $k + m + n$ can be left out.

\section{Multiplication}

\begin{definition}[Multiplication]
\label{def:multiplication}
Let $m \in \N$, and let $g \colon \N \to \N$ be the map with $x \via{g} x + m$. By
\cref{thm:recursion} there is exactly one $h \colon \N \to \N$ with $0 \via{h} 0$
and $h \circ S = g \circ h$. \emph{Multiplication} is the operation on $\N$ with
$(m, n) \via{\cdot} k \defiff n \via{h} k$. Its value is written $m \cdot n$. It
is characterised by
\[
  \text{(M1)}\quad m \cdot 0 = 0, \qquad\qquad \text{(M2)}\quad m \cdot (n + 1) = m
  \cdot n + m.
\]
Products are evaluated before sums: $m \cdot n + m$ means $(m \cdot n) + m$.
\register{naturals}{$m \cdot n$}{product of $m$ and $n$; $\cdot$ is evaluated
before $+$}
\end{definition}

\begin{theorem}[Laws of multiplication]
\label{thm:multiplication-laws}
For all $k, m, n \in \N$:
\begin{enumerate}
  \item $0 \cdot n = 0$;
  \item $(m + 1) \cdot n = m \cdot n + n$;
  \item $m \cdot n = n \cdot m$ (commutativity);
  \item $k \cdot (m + n) = k \cdot m + k \cdot n$ and $(m + n) \cdot k = m \cdot k
    + n \cdot k$ (distributivity);
  \item $(k \cdot m) \cdot n = k \cdot (m \cdot n)$ (associativity);
  \item $1 \cdot n = n = n \cdot 1$ ($1$ is neutral);
  \item $m \cdot n = 0 \Rightarrow m = 0 \lor n = 0$.
\end{enumerate}
\end{theorem}

\begin{proof}
(1)--(5) by induction on $n$, using (M1), (M2) and \cref{thm:addition-laws}.

(1) $0 \cdot 0 = 0$. If $0 \cdot n = 0$, then $0 \cdot (n + 1) = 0 \cdot n + 0 =
0$.

(2) $(m + 1) \cdot 0 = 0 = m \cdot 0 + 0$. If it holds for $n$:
\begin{align*}
  (m + 1) \cdot (n + 1) &= (m + 1) \cdot n + (m + 1) = (m \cdot n + n) + (m + 1) \\
  &= m \cdot n + ((n + m) + 1) = m \cdot n + (m + (n + 1)) \\
  &= (m \cdot n + m) + (n + 1) = m \cdot (n + 1) + (n + 1).
\end{align*}

(3) $m \cdot 0 = 0 = 0 \cdot m$ by (1). If $m \cdot n = n \cdot m$, then $m \cdot
(n + 1) = m \cdot n + m = n \cdot m + m = (n + 1) \cdot m$ by (2).

(4) $k \cdot (m + 0) = k \cdot m = k \cdot m + k \cdot 0$. If it holds for $n$,
then
\[
  k \cdot (m + (n + 1)) = k \cdot ((m + n) + 1) = k \cdot (m + n) + k = (k \cdot m
  + k \cdot n) + k = k \cdot m + k \cdot (n + 1).
\]
The second form follows by (3).

(5) $(k \cdot m) \cdot 0 = 0 = k \cdot (m \cdot 0)$. If it holds for $n$, then
$(k \cdot m) \cdot (n + 1) = (k \cdot m) \cdot n + k \cdot m = k \cdot (m \cdot n)
+ k \cdot m = k \cdot (m \cdot n + m) = k \cdot (m \cdot (n + 1))$, using (4).

(6) $n \cdot 1 = n \cdot (0 + 1) = n \cdot 0 + n = n$, and $1 \cdot n = n$ by (3).

(7) If $m \neq 0$ and $n \neq 0$, write $m = p + 1$ and $n = q + 1$
(\cref{cor:predecessor}). Then $m \cdot n = m \cdot q + (p + 1) = (m \cdot q + p) +
1 \neq 0$.
\end{proof}

\section{Order}

\begin{definition}[Order on $\N$]
\label{def:order-naturals}
For $m, n \in \N$:
\[
  m < n \defiff m \in n, \qquad m \le n \defiff m < n \lor m = n.
\]
\register{naturals}{$m < n$, $m \le n$}{order of natural numbers: $m \in n$, and
$m \in n \lor m = n$}
\end{definition}

So $n = \setof{m \in \N}{m < n}$: a natural number is the set of the numbers
before it, as promised. (Elements of $n$ are natural numbers by
\cref{lem:shape}(1).)

\begin{lemma}
\label{lem:order-naturals}
For all $m, n \in \N$:
\begin{enumerate}
  \item $m < n + 1 \Leftrightarrow m \le n$;
  \item $\lnot (n < 0)$;
  \item $0 \le n$;
  \item $m < n \Rightarrow m + 1 \le n$.
\end{enumerate}
\end{lemma}

\begin{proof}
(1) $n + 1 = n \cup \set{n}$, so $m \in n + 1 \Leftrightarrow m \in n \lor m = n$.

(2) $0 = \varnothing$ has no elements.

(3) By induction: $0 \le 0$, and $0 \le n$ gives $0 < n + 1$ by (1).

(4) By induction on $n$. For $n = 0$ there is nothing to show, by (2). Assume it
holds for $n$ and let $m < n + 1$, so $m \le n$ by (1). If $m = n$, then $m + 1 = n
+ 1$. If $m < n$, then $m + 1 \le n$ by the induction hypothesis, so $m + 1 < n +
1$ by (1).
\end{proof}

\begin{theorem}[$\le$ is a total order]
\label{thm:total-order-naturals}
$<$ is irreflexive and transitive on $\N$, and \emph{trichotomy} holds: for all
$m, n$ exactly one of
\[
  m < n, \qquad m = n, \qquad n < m
\]
is true. Consequently $\le$ is a total order on $\N$ (\cref{def:order}).
\end{theorem}

\begin{proof}
Irreflexive: $n \notin n$ (\cref{lem:shape}(3)). Transitive: if $k \in m$ and $m
\in n$, then $m \subseteq n$ (\cref{lem:shape}(2)), so $k \in n$.

\emph{At least one} holds, by induction on $n$ for fixed $m$. For $n = 0$: $0 \le
m$ by \cref{lem:order-naturals}(3), so $0 = m$ or $0 < m$. Step: assume one of the
three holds for $n$. If $m < n$ or $m = n$, then $m \le n$, so $m < n + 1$ by
\cref{lem:order-naturals}(1). If $n < m$, then $n + 1 \le m$ by
\cref{lem:order-naturals}(4). Either way one of the three holds for $n + 1$.

\emph{At most one} holds: $m < n$ together with $m = n$ or with $n < m$ would give
$m < m$ by transitivity, against irreflexivity.

By \cref{lem:strict-order}, $\le$ is a partial order, and trichotomy makes it
total.
\end{proof}

\begin{theorem}[Order and arithmetic]
\label{thm:order-arithmetic}
For all $k, m, n \in \N$:
\begin{enumerate}
  \item $m \le n \Leftrightarrow \exists d \in \N\; n = m + d$, and this $d$ is
    unique;
  \item $m < n \Leftrightarrow \exists d \in \N\; n = m + (d + 1)$;
  \item $m < n \Rightarrow m + k < n + k$;
  \item $m < n \land k \neq 0 \Rightarrow m \cdot k < n \cdot k$.
\end{enumerate}
\end{theorem}

\begin{proof}
(1) ($\Leftarrow$) By induction on $d$: $m \le m + 0$. If $m \le m + d$, then $m +
d < (m + d) + 1 = m + (d + 1)$ (\cref{lem:order-naturals}(1)), so $m < m + (d + 1)$
by transitivity. ($\Rightarrow$) By induction on $n$. If $m \le 0$, then $m = 0$
(\cref{lem:order-naturals}(2)), and $d = 0$ works. Assume the claim for $n$, and
let $m \le n + 1$. If $m = n + 1$, take $d = 0$. Otherwise $m < n + 1$, so $m \le
n$, so $n = m + d$ for some $d$, and $n + 1 = (m + d) + 1 = m + (d + 1)$.
Uniqueness: $m + d = m + d'$ gives $d = d'$ by \cref{thm:addition-laws}(3) and (5).

(2) ($\Rightarrow$) $m + 1 \le n$ (\cref{lem:order-naturals}(4)), so by (1) $n = (m
+ 1) + d = m + (d + 1)$, using \cref{thm:addition-laws}. ($\Leftarrow$) $m \le m +
d < (m + d) + 1 = n$ by (1) and \cref{lem:order-naturals}(1).

(3) By (2), $n = m + (d + 1)$, so $n + k = (m + k) + (d + 1)$ by
\cref{thm:addition-laws}(3) and (4). Use (2) again.

(4) By (2), $n = m + (d + 1)$, so $n \cdot k = m \cdot k + (d + 1) \cdot k$. Since
$k \neq 0$ and $d + 1 \neq 0$, $(d + 1) \cdot k \neq 0$
(\cref{thm:multiplication-laws}(7)), so $(d + 1) \cdot k = e + 1$ for some $e$
(\cref{cor:predecessor}). Use (2).
\end{proof}

\begin{corollary}[Cancellation]
\label{cor:cancellation-naturals}
$m + k < n + k \Rightarrow m < n$. And if $k \neq 0$: $m \cdot k = n \cdot k
\Rightarrow m = n$, and $m \cdot k < n \cdot k \Rightarrow m < n$.
\end{corollary}

\begin{proof}
Each time, compare $m$ and $n$ by trichotomy (\cref{thm:total-order-naturals}).
The cases other than the claimed one contradict the hypothesis by
\cref{thm:order-arithmetic}(3) and (4). For example, if $m \cdot k = n \cdot k$ but
$m < n$, then $m \cdot k < n \cdot k$, contradicting irreflexivity.
\end{proof}

\section{Well-ordering}

\begin{theorem}[Well-ordering principle]
\label{thm:well-ordering}
Every nonempty subset of $\N$ has a least element. So $(\N, \le)$ is well-ordered
(\cref{def:well-order}).
\end{theorem}

\begin{proof}
Let $X \subseteq \N$ have no least element. We show $X = \varnothing$. Let
\[
  Y \coloneqq \setof{n \in \N}{\forall m \in \N\; (m < n \Rightarrow m \notin X)}
\]
be the set of $n$ such that no number below $n$ is in $X$. We show $Y$ is
inductive. $0 \in Y$ vacuously (\cref{lem:order-naturals}(2)). Let $n \in Y$. If
$n$ were in $X$, it would be the least element of $X$: for $x \in X$, $x < n$ is
impossible because $n \in Y$, so $n \le x$ by trichotomy. So $n \notin X$. Then no
$m \le n$ is in $X$, so no $m < n + 1$ (\cref{lem:order-naturals}(1)). So $n + 1
\in Y$.

By induction $Y = \N$. Every $n \in \N$ satisfies $n < n + 1$ and $n + 1 \in Y$, so
$n \notin X$.
\end{proof}

\begin{corollary}[Strong induction]
\label{cor:strong-induction}
Let $P$ be a predicate such that for every $n \in \N$,
\[
  \bigl( \forall m \in \N\; (m < n \Rightarrow P(m)) \bigr) \Rightarrow P(n).
\]
Then $P(n)$ holds for every $n \in \N$.
\end{corollary}

\begin{proof}
Suppose $X \coloneqq \setof{n \in \N}{\lnot P(n)}$ is nonempty, and let $n$ be its
least element (\cref{thm:well-ordering}). Every $m < n$ is outside $X$, so $P(m)$
holds for all $m < n$. The hypothesis gives $P(n)$, contradicting $n \in X$.
\end{proof}

The case $n = 0$ has no $m < 0$, so it requires proving $P(0)$ outright. Strong
induction lets the step use \emph{all} smaller cases, not only the previous one. It
will be used for prime factorization.

\section{What is still missing}

$\N$ has addition and multiplication with all the familiar laws, and an order
compatible with both. But the equation $m + x = n$ only has a solution $x \in \N$
when $m \le n$ (\cref{thm:order-arithmetic}(1)). There is no $x$ with $1 + x = 0$.
\Cref{ch:integers} builds a larger number system where subtraction always works.

\begin{exercises}
  \item Prove $2 \cdot 3 = 6$ from the definitions, citing each step.
  \item Prove that there is no $m$ with $n < m$ and $m < n + 1$.
  \item Use \cref{thm:recursion} to define, for each $m$, the map $p \colon \N \to
    \N$ with $0 \via{p} 1$ and $n \via{p} y \Rightarrow n + 1 \via{p} y \cdot m$.
    Its value at $n$ is the $n$-th power of $m$. Show that the value at $a + b$ is
    the product of the values at $a$ and at $b$.
  \item Use \cref{cor:recursion-index} to define the \emph{factorial}, the map $F
    \colon \N \to \N$ with $0 \via{F} 1$ and $n \via{F} y \Rightarrow n + 1 \via{F}
    y \cdot (n + 1)$. Find the value of $F$ at $4$.
  \item For a sequence $x \colon \N \to \N$, use \cref{cor:recursion-index} to
    define its \emph{partial sums}, the sequence $s$ with $0 \via{s} 0$ and $n
    \via{s} y \land n \via{x} z \Rightarrow n + 1 \via{s} y + z$. Let $x = \id_\N$.
    Prove: if $n + 1 \via{s} y$, then $2 \cdot y = n \cdot (n + 1)$.
  \item Show that $\N$ has no greatest element.
\end{exercises}
